% ============================================================
\chapter{Th\'eor\`emes de Comparaison}
\label{ch:comparaison}
% ============================================================

\section{Introduction et motivation}

L'une des id\'ees les plus f\'econdes de la g\'eom\'etrie riemannienne repose sur un principe remarquablement simple~: si l'on sait que la courbure d'une vari\'et\'e est born\'ee, alors sa g\'eom\'etrie ne peut pas s'\'ecarter trop de celle d'un espace \og mod\`ele \fg{} de courbure constante. Ce programme de comparaison, initi\'e par Harry Rauch dans les ann\'ees 1950 puis d\'evelopp\'e par Marcel Berger, Wilhelm Klingenberg, Jeff Cheeger et Detlef Gromoll dans les d\'ecennies suivantes, a produit certains des r\'esultats les plus profonds de la g\'eom\'etrie globale. Le th\'eor\`eme de Rauch compare le comportement des champs de Jacobi~; celui de Bishop--Gromov compare les volumes~; celui de Toponogov compare les triangles. Chacun extrait des conclusions g\'eom\'etriques globales \`a partir d'une hypoth\`ese locale sur la courbure, illustrant la puissance de la courbure comme invariant de contr\^ole.

\section{\'Equation de Jacobi et solutions mod\`eles}

Rappelons que le long d'une g\'eod\'esique $\gamma$ param\'etr\'ee par longueur d'arc,
un champ de Jacobi orthogonal $J \perp \dot\gamma$ satisfait~:
\[
    J'' + K(t)\, J = 0,
\]
o\`u $K(t)$ repr\'esente la courbure sectionnelle du plan $(\dot\gamma, J)$.

\begin{definition}[Solutions mod\`eles]
Pour $\kappa \in \R$, la solution de $f'' + \kappa f = 0$ avec $f(0) = 0$,
$f'(0) = 1$ est~:
\[
    \operatorname{sn}_\kappa(t) = \begin{cases}
    \frac{1}{\sqrt{\kappa}}\sin(\sqrt{\kappa}\, t) & \text{si } \kappa > 0, \\
    t & \text{si } \kappa = 0, \\
    \frac{1}{\sqrt{-\kappa}}\sinh(\sqrt{-\kappa}\, t) & \text{si } \kappa < 0.
    \end{cases}
\]
\end{definition}

\begin{formules}[title=\textbf{Fonctions de comparaison}]
\begin{align*}
    \operatorname{sn}_\kappa(t) &= \text{solution de } f'' + \kappa f = 0, \; f(0)=0, \; f'(0)=1, \\
    \operatorname{cn}_\kappa(t) &= \operatorname{sn}_\kappa'(t) = \text{solution de }
    f'' + \kappa f = 0, \; f(0)=1, \; f'(0)=0.
\end{align*}
\end{formules}

\section{Lemme de comparaison de Sturm}

\begin{lemme}[Sturm]
\label{lem:sturm}
Soient $f$ et $g$ deux solutions des \'equations~:
\[
    f'' + K_1(t)\, f = 0, \qquad g'' + K_2(t)\, g = 0,
\]
avec $f(0) = g(0) = 0$, $f'(0) = g'(0) = 1$, et $K_1(t) \geq K_2(t)$.

Alors $f(t) \leq g(t)$ sur l'intervalle $[0, t_1]$ o\`u $t_1$ est le premier z\'ero
de $f$.
\end{lemme}

\begin{proof}
On consid\`ere la fonction $h = f'g - fg'$. On a $h(0) = 0$ et~:
\[
    h' = f''g - fg'' = -K_1 fg + K_2 fg = (K_2 - K_1)fg \leq 0,
\]
tant que $f \geq 0$ et $g \geq 0$. Donc $h \leq 0$, i.e.\ $f'/f \leq g'/g$
(quand $f, g > 0$). Int\'egrant~: $\ln f - \ln g$ est d\'ecroissante, et comme
$\lim_{t \to 0^+} f(t)/g(t) = 1$, on obtient $f \leq g$.
\end{proof}

\section{Th\'eor\`eme de comparaison de Rauch}

\begin{theoreme}[Rauch]
\label{thm:rauch}
Soient $(M, g)$ et $(\bar{M}, \bar{g})$ deux vari\'et\'es riemanniennes, et
$\gamma : [0, L] \to M$, $\bar\gamma : [0, L] \to \bar{M}$ deux g\'eod\'esiques
param\'etr\'ees par longueur d'arc, sans points conjugu\'es. Supposons que~:
\[
    K_M(\dot\gamma, \cdot) \leq K_{\bar{M}}(\dot{\bar\gamma}, \cdot).
\]
Si $J$ et $\bar{J}$ sont des champs de Jacobi orthogonaux le long de $\gamma$
et $\bar\gamma$ avec $J(0) = \bar{J}(0) = 0$ et $\norm{J'(0)} = \norm{\bar{J}'(0)}$,
alors~:
\[
    \norm{J(t)} \geq \norm{\bar{J}(t)}, \quad \forall\, t \in [0, L].
\]
\end{theoreme}

\begin{proof}[Id\'ee de la preuve]
On consid\`ere le quotient $u(t) = \norm{J(t)}/\operatorname{sn}_\kappa(t)$
o\`u $\kappa$ est la borne inf\'erieure de la courbure. Par le lemme de Sturm,
on compare $\norm{J}$ avec la solution mod\`ele. La cl\'e est l'in\'egalit\'e~:
\[
    \frac{d}{dt}\frac{\norm{J}'}{\norm{J}} \leq -K(t),
\]
qui se d\'eduit de l'\'equation de Jacobi et de l'in\'egalit\'e de Cauchy--Schwarz.
On conclut par le th\'eor\`eme de comparaison de Sturm.
\end{proof}

\begin{corollaire}[Comparaison avec l'espace mod\`ele]
Si $K_M \leq \kappa$ (borne sup\'erieure), alors~:
\[
    \norm{J(t)} \geq \operatorname{sn}_\kappa(t)\, \norm{J'(0)}.
\]
Si $K_M \geq \kappa$ (borne inf\'erieure), alors~:
\[
    \norm{J(t)} \leq \operatorname{sn}_\kappa(t)\, \norm{J'(0)}.
\]
\end{corollaire}

\section{Th\'eor\`eme de Toponogov}

\begin{theoreme}[Toponogov -- version globale]
\label{thm:toponogov}
Soit $(M, g)$ une vari\'et\'e riemannienne compl\`ete avec $K \geq \kappa$. Soit
$\triangle(p, q, r)$ un triangle g\'eod\'esique dans $M$ (form\'e de trois g\'eod\'esiques
minimisantes). Soit $\bar\triangle(\bar{p}, \bar{q}, \bar{r})$ un triangle de
comparaison dans $M_\kappa$ (m\^emes longueurs de c\^ot\'es). Alors~:

\textbf{Version angles~:} Les angles du triangle dans $M$ sont \emph{sup\'erieurs
ou \'egaux} aux angles correspondants du triangle de comparaison~:
\[
    \angle_p(q, r) \geq \bar\angle_{\bar{p}}(\bar{q}, \bar{r}).
\]

\textbf{Version distances~:} Pour tout point $x$ sur le c\^ot\'e $[q, r]$ et le
point correspondant $\bar{x}$ sur $[\bar{q}, \bar{r}]$ (m\^emes proportions)~:
\[
    d(p, x) \leq d(\bar{p}, \bar{x}).
\]
\end{theoreme}

\begin{attention}[title=\textbf{Conditions d'application de Toponogov}]
Le th\'eor\`eme de Toponogov requiert~:
\begin{enumerate}
    \item La compl\'etude de $(M, g)$.
    \item Une borne \emph{inf\'erieure} sur la courbure sectionnelle.
    \item Les c\^ot\'es du triangle sont des g\'eod\'esiques \emph{minimisantes}.
    \item Si $\kappa > 0$, le p\'erim\`etre du triangle ne d\'epasse pas $2\pi/\sqrt{\kappa}$.
\end{enumerate}
\end{attention}

\begin{corollaire}
Si $K \geq 0$ (courbure non n\'egative), les triangles g\'eod\'esiques sont
<<~plus gros~>> que les triangles euclidiens correspondants~: les angles sont plus
grands et les <<~m\'edianes~>> plus courtes.
\end{corollaire}

\section{Th\'eor\`eme de comparaison de Bishop--Gromov}

\begin{theoreme}[Bishop--Gromov]
\label{thm:bishop_gromov}
Soit $(M^n, g)$ une vari\'et\'e riemannienne compl\`ete avec
$\operatorname{Ric} \geq (n-1)\kappa$. Soit $V_\kappa(r)$ le volume de la boule
de rayon $r$ dans l'espace mod\`ele $M^n_\kappa$. Alors la fonction~:
\[
    r \mapsto \frac{\operatorname{Vol}(B(p, r))}{V_\kappa(r)}
\]
est d\'ecroissante pour $r > 0$, et tend vers $1$ quand $r \to 0$.

En particulier~:
\[
    \operatorname{Vol}(B(p, r)) \leq V_\kappa(r), \quad \forall\, r > 0.
\]
\end{theoreme}

\begin{proof}[Id\'ee de la preuve]
En coordonn\'ees polaires g\'eod\'esiques~:
$\operatorname{Vol}(B(p,r)) = \int_{S^{n-1}} \int_0^{\min(r, c(v))}
\mathcal{A}(t, v)\, dt\, d\sigma(v)$,
o\`u $\mathcal{A}(t, v)$ est le jacobien de $\exp_p$ et $c(v)$ la distance au point
de coupure. Par l'\'equation de Jacobi et la borne $\operatorname{Ric} \geq (n-1)\kappa$,
on montre que $\mathcal{A}(t,v)/\mathcal{A}_\kappa(t)$ est d\'ecroissante en $t$,
o\`u $\mathcal{A}_\kappa(t) = \operatorname{sn}_\kappa(t)^{n-1}$ est le jacobien
de l'espace mod\`ele.
\end{proof}

\begin{corollaire}[Comparaison de volumes]
Sous les hypoth\`eses du th\'eor\`eme, pour $0 < r \leq R$~:
\[
    \frac{\operatorname{Vol}(B(p, R))}{\operatorname{Vol}(B(p, r))}
    \leq \frac{V_\kappa(R)}{V_\kappa(r)}.
\]
\end{corollaire}

\begin{corollaire}[Croissance du volume]
Si $\operatorname{Ric} \geq 0$, alors
$\operatorname{Vol}(B(p, r)) \leq \omega_n r^n$ (volume euclidien).
\end{corollaire}

\section{Application~: th\'eor\`eme de Cheng}

\begin{theoreme}[Cheng]
Soit $(M^n, g)$ compl\`ete avec $\operatorname{Ric} \geq (n-1)\kappa > 0$.
Si $\operatorname{diam}(M) = \pi/\sqrt{\kappa}$ (\'egalit\'e dans Bonnet--Myers),
alors $M$ est isom\'etrique \`a $S^n(1/\sqrt{\kappa})$.
\end{theoreme}

\section{Th\'eor\`eme de comparaison de Laplacien}

\begin{theoreme}[Comparaison du laplacien]
Soit $(M^n, g)$ avec $\operatorname{Ric} \geq (n-1)\kappa$. La fonction distance
$r(x) = d(p, x)$ v\'erifie, au sens des distributions~:
\[
    \Delta r \leq (n-1)\frac{\operatorname{cn}_\kappa(r)}{\operatorname{sn}_\kappa(r)}.
\]
Pour $\kappa = 0$~: $\Delta r \leq \frac{n-1}{r}$.
Pour $\kappa = 1$~: $\Delta r \leq (n-1)\cot r$.
Pour $\kappa = -1$~: $\Delta r \leq (n-1)\coth r$.
\end{theoreme}

\section{Th\'eor\`eme de Gromov sur le nombre de g\'en\'erateurs}

\begin{theoreme}[Gromov]
Si $(M^n, g)$ est compl\`ete avec $\operatorname{Ric} \geq 0$, alors le groupe
fondamental $\pi_1(M)$ est engendr\'e par au plus $C(n)$ \'el\'ements, o\`u $C(n)$
ne d\'epend que de la dimension.
\end{theoreme}

\section{Espace de Gromov--Hausdorff}

\begin{definition}[Distance de Gromov--Hausdorff]
La \emph{distance de Gromov--Hausdorff} entre deux espaces m\'etriques compacts
$(X, d_X)$ et $(Y, d_Y)$ est~:
\[
    d_{GH}(X, Y) = \inf \{ d_H(\varphi(X), \psi(Y)) \},
\]
o\`u l'infimum porte sur toutes les isom\'etries $\varphi : X \to Z$,
$\psi : Y \to Z$ dans un espace m\'etrique $(Z, d_Z)$, et $d_H$ est la distance
de Hausdorff.
\end{definition}

\begin{theoreme}[Pr\'e-compacit\'e de Gromov]
L'ensemble des vari\'et\'es riemanniennes $(M^n, g)$ v\'erifiant
$\operatorname{Ric} \geq (n-1)\kappa$ et $\operatorname{diam}(M) \leq D$
est pr\'e-compact pour la distance de Gromov--Hausdorff.
\end{theoreme}

\section{Exercices}

\begin{exercice}
V\'erifier le th\'eor\`eme de Bishop--Gromov pour les boules de $S^n$ de courbure $1$,
en calculant explicitement $\operatorname{Vol}(B(p, r))$ et en comparant avec
$V_1(r)$.
\end{exercice}

\begin{exercice}
Montrer que si $(M, g)$ est compl\`ete avec $K \geq 1$, alors
$\operatorname{diam}(M) \leq \pi$.
\end{exercice}

\begin{exercice}
Appliquer le th\'eor\`eme de Toponogov pour montrer que si $K \geq 0$ et $M$ est
non compacte, alors $M$ contient une droite g\'eod\'esique (th\'eor\`eme de la droite).
\end{exercice}

\begin{exercice}
Utiliser la comparaison de volumes pour montrer que si $\operatorname{Ric} \geq 0$
et $\operatorname{Vol}(B(p,r)) \geq c\, r^n$ pour un $c > 0$ et tout $r$, alors
$M$ a au plus un nombre fini de bouts.
\end{exercice}

\begin{exercice}
Calculer la comparaison de Laplacien pour la fonction distance sur $\mathbb{H}^3$
et v\'erifier qu'on obtient $\Delta r = 2\coth r$.
\end{exercice}

\begin{exercice}
Montrer que si $K \leq 0$ et $M$ est simplement connexe, alors il n'y a pas
de points conjugu\'es (en utilisant le th\'eor\`eme de comparaison de Rauch).
\end{exercice}
